\documentclass[11pt]{article}
\usepackage[margin=1in]{geometry}
\usepackage{enumitem}
\usepackage{amsmath}
\usepackage{hyperref}
\hypersetup{colorlinks=true, urlcolor=blue, linkcolor=blue}
\usepackage{parskip}

\begin{document}

\begin{center}
  {\Large\bfseries Dengue Exercise (Part 2): From Least Squares to ARGO}\\[2pt]
  {\normalsize SISMID 2026 \textbullet\ Day 2, 1:30 \textbullet\ From the Dengue Exercise to ARGO}
\end{center}

\medskip

In the previous exercise you fit a \emph{static} model to predict dengue cases in Mexico
using training data from 2004--2006: a single, unchanging line you use to predict all future
cases. In this exercise you fit a \emph{dynamic} model that is retrained at each prediction
step, then add more features until the model becomes a basic \textbf{ARGO} model
(AutoRegression with Google search data). A short final part sketches the network
(\textbf{ARGONet}) idea.

You may work in \textbf{Python or R}, and either drive a coding agent (Codex, Claude Code,
or Antigravity CLI) with a prompt per part, or fill in the pre-filled notebook. Data:
\texttt{MX\_Dengue\_trends.csv} (\texttt{Dengue CDC} = reported cases; \texttt{dengue},
\texttt{sintomas de dengue}, \texttt{mosquito}, \texttt{dengue sintomas} = Google search
interest), from \url{https://github.com/MIGHTE-lab/SISMID25}.

\begin{enumerate}[label=(\alph*)]
  \item \textbf{Dynamic (sliding-window) training.} In a sliding-window approach, the
        original training period (Jan.\ 2004 -- Dec.\ 2006) fits a least squares model and
        predicts the next time step (Jan.\ 2007). For the next prediction, the training
        period shifts by one month (Feb.\ 2004 -- Jan.\ 2007) and the model is retrained.
        Continue this for each prediction from 2007--2011 and store the predictions in a
        vector.

  \item \textbf{Compare to static.} Plot the cases predicted by the dynamic method against
        the predictions from the static model and the ground truth.

  \item \textbf{Score both.} Compute the mean square error (MSE, or RMSE) of both methods
        against the ground truth. Which is better, and why?

  \item \textbf{Add covariates.} Instead of just \texttt{dengue}, add the remaining Google
        search terms as predictors in the dynamic model. How does the error change? With
        many correlated terms, try \textbf{Lasso} ($\ell_1$-penalized regression) so useless
        terms drop out on their own.

  \item \textbf{Add autoregression (this is ARGO).} Use the fact that the model is refit each
        step to add past observations of the ground truth as features. For the Jan.\ 2007
        prediction, add the observed case counts from Nov.\ 2006 and Dec.\ 2006 (a
        Google\,+\,AR2 model). Implement this dynamic model with autoregression and compare
        its error to the others. This combination,
        \[
          y_t \;=\; \mu \;+\; \sum_{k=1}^{2}\beta_k\,y_{t-k}
            \;+\; \sum_{j}\alpha_j\,x_{j,t} \;+\; \varepsilon_t,
        \]
        trained dynamically, is a basic \textbf{ARGO} model.

        \emph{Hint:} shift the case-count column and stack it onto the feature matrix; the
        first one or two windows may have fewer than 36 observations.

  \item \textbf{Discuss and extend (ARGONet).} You should see the error fall at each step:
        static $\to$ dynamic $\to$ +terms $\to$ +autoregression. Now suppose you had
        \emph{several} locations (states or counties) instead of one country. \textbf{ARGONet}
        adds a network step: borrow signal from connected neighbors, for example by adding a
        neighbor's recent activity (or its own ARGO estimate) as an extra predictor, choosing
        neighbors by shared movement or short-lag correlation. Sketch, in words or code, how
        you would add one neighbor's lagged signal to the model above, and say when you would
        expect it to help.
\end{enumerate}

\vfill
\noindent\textit{Stretch:} refit ARGO with an $\ell_1$ penalty and a two-year rolling window
(the full ARGO of Yang, Santillana \& Kou, PNAS 2015), and compare against the plain
least-squares version.

\end{document}
